% ======================================================================
% SECTION 4: HR 9504 - Physical AI Clinical Error Reduction Act of 2026
% Adaption of HTI-1 DSI Final Rule 2023 and FDA AI Decision-Making
% Draft Guidance 2025.
%
% This file consolidates instructions previously housed in the prior
% template's Patient Priority Framework dimension 2.6 (AI-exclusion
% contestability), Implementation/Metrics (Domain 3 control-in-action
% and Domain 4 fairness/technical quality), Discussion (11.1.e ONC
% HTI-1 extension, 11.2.b ePRO evidence, 11.2.d Mazor null result,
% 11.3.c FDA AI Draft 2025), Limitations (12.1.d HR 4504 limitation),
% and Conclusions (13.1 fourth list item) into the body of HR 9504.
% ======================================================================

[CONSOLIDATED PATIENT PRIORITY FRAMEWORK DIMENSION 2.6 (formerly
Patient Priority Framework subsection 2.6). AI-exclusion
contestability: HTI-1 DSI transparency, FDA AI draft guidance 2025,
mandatory appeal rights for any patient excluded from a trial or
treatment by an AI gatekeeping decision. Cite
\cite{rank9_onc_hti1_2023}, \cite{hti1_dsi_fact_sheet_2023}, and
\cite{fda_ai_regulatory_decision_draft_2025}.]

[CONSOLIDATED INTRODUCTION TIE-IN (formerly Introduction block 2 item
2.f). ONC HTI-1 Final Rule (2023) and HTI-1 DSI Fact Sheet impose
algorithm-transparency requirements for predictive decision support
interventions in certified health IT. HR 9504 extends transparency
into mandatory error-rate tracking and patient-initiated AI
second-opinion rights. Cite \cite{rank9_onc_hti1_2023} and
\cite{hti1_dsi_fact_sheet_2023}.]

[PROCESSING INSTRUCTION (HR 9504 - Physical AI Clinical Error
Reduction Act of 2026, populate to 10 to 13 paragraphs structured
into the five legislative-act subsections, plus one ASCII comparison
diagram):

The bill must explicitly state in its first paragraph that it is an
"Adaption of HTI-1 DSI Final Rule of 2023 and FDA AI Regulatory
Decision-Making Draft Guidance 2025" so that readers can trace the
chain of prior authority. Cite \cite{rank9_onc_hti1_2023},
\cite{hti1_dsi_fact_sheet_2023}, and
\cite{fda_ai_regulatory_decision_draft_2025}.

Subsection 4.1 - Findings and Declarations: This is the core bill
for the project brief's third aim ("reduce human doctor and nurse
error rates and lack of physical and intellectual capabilities").
Read the following in full:
- new-trial/national-24-7-trial/paper/full-paper/sections/results.tex
  (the four LLM simulations demonstrating that AI patient prediction
  outperforms current narrow supervised models)
- new-trial/national-24-7-trial/Background-A/chunk_01_baseline_and_short_horizon.md
  and chunk_02_multimodal_and_limitations.md (Manz 2020 AUC 0.89,
  Sidey-Gibbons 2022, the Huang 2025 null result of ML versus Cox)
- new-trial/national-24-7-trial/Background-B/chunk_01_baseline_and_prediction_domains.md
  and chunk_02_response_metrics_conclusions.md (SCORPIO, PROGPATH,
  multimodal foundation models)
- patients/paper/Deep-Research-3/part2_ai_robotics_legislation.md
  (the SHIELD-RT prospective RCT: 45 percent reduction in acute-care
  events, 48 percent reduction in costs)

Compose 6 to 8 enumerated findings citing
\cite{kawchak_2026_19994945} for the four-simulation paper,
\cite{vorontsov2024virchow}, \cite{rocque2025remote_symptom_monitoring},
\cite{gonzalez2025bayesian_urgent_care_nsclc},
\cite{mazor2025ai_notifications_trial}, and the FDA RTCT announcement
at \cite{fda_rtct_press_2026}. The findings must:
- Acknowledge respectfully that human doctors and nurses make errors
  at documented rates (cite the prior published rates rather than
  inflating them).
- Acknowledge that AI and robotics can reduce specific error
  categories where prospectively validated (SHIELD-RT, ASyMS,
  eRAPID).
- Acknowledge that humans retain irreplaceable roles in empathy,
  consent counseling, and ethical oversight, and that the bill does
  NOT eliminate human doctors or nurses; it shifts the patient's
  default for procedurally definable tasks to physical-AI execution
  where the patient prefers and the evidence supports.

Subsection 4.2 - Definitions:
- "Documented human error rate" means the per-task error rate
  published in peer-reviewed literature for the specific task
  category (surgical site infection, medication reconciliation,
  needle placement deviation, RT positioning, sepsis recognition,
  etc.).
- "Physical AI substitute" means a Physical AI system, robot, or
  agent that can perform the same task category with a documented
  error rate at or below the human baseline. Reference the 53 core
  agents at sponsor/final_paper/scripts/core_agents/ (in particular
  safety_agent.py, clinical_accountability_agent.py, and
  audit_trail_manager.py) for the operational baseline.
- "Patient-initiated AI second opinion" means the patient's right to
  request, before any consequential clinical decision, a parallel
  evaluation by a qualified AI system meeting HTI-1 DSI transparency
  requirements.

Subsection 4.3 - Error Rate Reduction Rights (operative core):
- Right to a published per-site, per-task human error rate prior to
  consenting to any procedure.
- Right to substitute a Physical AI executor for any task where the
  documented Physical AI error rate is below the human baseline AND
  the patient prefers AI execution.
- Right to a patient-initiated AI second opinion, generated by a
  qualified system, prior to any consequential clinical decision
  (surgery, chemotherapy regimen change, end-of-life care decision).
- Right to mandatory error tracking across all sites, with public
  per-site reporting; reference the dashboard at
  sponsor/final_paper/scripts/dashboard/ and the 168-hour cumulative
  metrics at sponsor/final_paper/168_hours/instructions/core_i5_6200u_4gb/sponsor_168h_summary.json
  for the operational reporting pattern.
- Right to refuse a human-only execution path for any task where a
  qualified Physical AI substitute exists, mirroring the FDA RTCT
  pilot's "real-time signal sharing" model. Cite
  \cite{fda_rtct_press_2026}.
- Right to AI-exclusion contestability: any patient excluded from a
  trial, treatment, or procedural pathway by an automated decision
  may appeal under the HTI-1 DSI transparency framework. Cite
  \cite{rank9_onc_hti1_2023} and \cite{hti1_dsi_fact_sheet_2023}.

After Subsection 4.3, embed a comprehensive ASCII comparison diagram
in a verbatim block. Three columns (Task Category, Human Baseline
Error Rate, Physical AI Error Rate) and rows for: surgical site
infection, medication reconciliation, needle placement deviation,
RT positioning, sepsis recognition, AE detection from ePRO. Source
the human baselines from the prior literature cited in
patients/paper/Deep-Research-1/chunk_1_foundations.md and the AI
baselines from
new-trial/national-24-7-trial/Background-A/ and
sponsor/final_paper/scripts/output/sponsor_24h_summary.json.

Subsection 4.4 - Implementation: 18-month deadline for ALL
CMS-billable sites to publish per-task human and AI error rates;
HTI-1 DSI transparency requirements apply to every AI second-opinion
system; FDA, AHRQ, ONC, and CMS roles. Reference the patient
priority operational baseline at
new-trial/national-24-7-trial/paper/full-paper/sections/discussion.tex.

Subsection 4.5 - Reporting and Enforcement: Annual federal
error-rate report by site, by task, by AI substitute (if any); CMS
reimbursement penalty for sites that fail to publish error rates;
information-blocking penalty harmonized with
\cite{federal_register_info_blocking_disincentives_2024}.]

[CONSOLIDATED IMPLEMENTATION-METRICS DOMAINS 3 AND 4 (formerly
Implementation Metrics subsection 10.2 Domains 3 and 4). The federal
Patient Control Dashboard must publish, for HR 9504, the following
metrics:

Domain 3 - Control-in-action: hospitalization rate, urgent-care or
ED use, alert response time, false-alert burden, serious adverse
events in home or robotics settings. Cite
\cite{rocque2025remote_symptom_monitoring},
\cite{gonzalez2025bayesian_urgent_care_nsclc},
\cite{dronca2026cancer_care_beyond_walls}, and
\cite{target2025robotic_bronchoscopy}.

Domain 4 - Fairness and technical quality: declared context of use,
external validation, subgroup performance, calibration monitoring,
update-control disclosure, revalidation after model changes. Cite
\cite{fda_ai_regulatory_decision_draft_2025} and
\cite{vorontsov2024virchow}.]

[CONSOLIDATED IMPLEMENTATION GUARDRAILS (formerly Implementation
Metrics subsection 10.3 guardrails (i), (ii), (iii) attached to HR
9504 because they all reference error reduction or AI substitution).

(i) No patient denied trial consideration by automated output alone.

(ii) Capability-neutral statutory rights that follow any AI
     architecture (LLM, foundation model, agentic system, robotics
     platform), drawing on
     \cite{fda_ai_regulatory_decision_draft_2025}.

(iii) Mandatory annual public reporting, auditability, and sunset
      review by Congress under FDORA Section 3209 authorities. Cite
      \cite{rank4_fdora_fda_modernization_2022}.]

[CONSOLIDATED DISCUSSION COMPARISON (formerly Discussion items
11.1.e ONC HTI-1 extension, 11.2.b ePRO evidence, 11.2.d Mazor null
result, 11.3.c FDA AI Draft 2025). HR 9504 extends ONC HTI-1 DSI
transparency into mandatory error-rate publication and
patient-initiated AI second-opinion rights. The ePRO and remote
symptom monitoring evidence (Rocque 2025, ASyMS, eRAPID) supports
HR 9504 without overclaiming. The Mazor 2025 null result on AI
notifications alone is critical and informs the bill's "no
automated-only denial" guardrail. Cite \cite{rank9_onc_hti1_2023},
\cite{rocque2025remote_symptom_monitoring}, \cite{kearney2009asyms},
\cite{maguire2017esmart}, \cite{absolom2021erapid}, and
\cite{mazor2025ai_notifications_trial}.]

[CONSOLIDATED LIMITATIONS PARAGRAPH (formerly Limitations item
12.1.d). HR 9504 is the most contentious of the seven because it
touches workforce questions beyond the project brief. The Huang
2025 null result indicates that ML does not always exceed Cox
regression on real-world structured survival data; the SHIELD-RT
prospective RCT is convincing for one domain (radiotherapy AE
prediction) but does not generalize without prospective validation
in each domain. Cite the four-simulation paper
\cite{kawchak_2026_19994945} and acknowledge the TRIPOD+AI / CREMLS
reporting gap. The bill respects human doctor and nurse roles in
empathy, consent counseling, and ethical oversight.]

[CONSOLIDATED CONCLUSIONS RESTATEMENT (formerly Conclusions item
13.1 fourth list item). HR 9504 constrains the AI-substitution
path to procedurally definable tasks with prospectively validated
error reductions, while explicitly preserving human roles in
empathy, consent counseling, and ethical oversight. Cite
\cite{kawchak_2026_19994945} and \cite{kawchak_2026_19119939}.]

Close the bill (1 paragraph after Subsection 4.5) with two patient-
control framings:

(i) FOR INDIVIDUAL PATIENTS: A 58-year-old female NSCLC patient
    matching PAT-2026-0042 reviews the published per-site error
    rates at her chosen site, requests an AI second opinion before
    her cycle-1 pembrolizumab dosing decision, and refuses a
    human-only RT positioning step in favor of the RTPOS-01 robot
    documented at hour-12 of
    new-trial/national-24-7-trial/hour-12/robot_logs.md.

(ii) FOR BROAD ADOPTION: All U.S. CMS-billable oncology trial sites
     must publish per-task error rates within 18 months. The bill
     positions the United States as Number 1 in the world for
     evidence-based patient-led error reduction. The bill explicitly
     respects the human doctor and nurse roles that remain
     irreplaceable (empathy, consent counseling, ethical oversight)
     and constrains the AI-substitution path to procedurally
     definable tasks with prospectively validated error reductions.
     Cite \cite{kawchak_2026_19994945} and
     \cite{kawchak_2026_19119939}.
